\newcommand{\IndexEntry}[2]{%
\noindent\parbox[t]{.89\linewidth}{\raggedright\hyperlink{#1}{#2}}%
\hfill\parbox[t]{.08\linewidth}{\raggedleft\scriptsize\hyperlink{#1}{\pageref{#1}}}%
\par\medskip}

\begin{frame}[label=nav-index-0-1]{Índice: Introducción}
\footnotesize\raggedright
\textcolor{structure.fg}{Subtemas y diapositivas (1/1)}\par\medskip
\begin{columns}[T,onlytextwidth]
\begin{column}{.48\textwidth}
\IndexEntry{nav-frame-001}{Introducción}
\IndexEntry{nav-frame-002}{Estrategia y alcance del recorrido}
\end{column}
\begin{column}{.48\textwidth}
\IndexEntry{nav-notation}{Convenciones de notación para todo el recorrido}
\end{column}
\end{columns}
\vfill
\end{frame}

\begin{frame}[label=nav-index-1-1]{Índice: Derivada de Lie e identidad principal (LL)}
\footnotesize\raggedright
\textcolor{structure.fg}{Subtemas y diapositivas (1/2)}\par\medskip
\begin{columns}[T,onlytextwidth]
\begin{column}{.48\textwidth}
\IndexEntry{nav-map-1-0}{Mapa de apertura}
\IndexEntry{nav-frame-004}{Derivada de Lie de un tensor general}
\IndexEntry{nav-frame-005}{Primer cálculo de $\Lie_\xi L$: usando que $L$ es escalar}
\IndexEntry{nav-frame-006}{Primer cálculo: compatibilidad métrica}
\IndexEntry{nav-frame-007}{Segundo cálculo de $\mathcal L_\xi L$: variando sus argumentos}
\IndexEntry{nav-frame-008}{Por qué aparece la derivada de Lie}
\end{column}
\begin{column}{.48\textwidth}
\IndexEntry{nav-frame-009}{Derivada de Lie del tensor de Riemann}
\IndexEntry{nav-frame-010}{Segundo cálculo: contribución métrica}
\IndexEntry{nav-frame-011}{Simetrías de \(P^{ijkl}\): primer par de índices}
\IndexEntry{nav-frame-012}{Simetrías de \(P^{ijkl}\): segundo par e intercambio}
\IndexEntry{nav-frame-013}{Por qué los cuatro términos son iguales}
\end{column}
\end{columns}
\vfill
\hyperlink{nav-index-1-2}{\scriptsize Más subtemas}\qquad
\hyperlink{nav-map-1-0}{\scriptsize Inicio de sección}
\end{frame}

\begin{frame}[label=nav-index-1-2]{Índice: Derivada de Lie e identidad principal (LL)}
\footnotesize\raggedright
\textcolor{structure.fg}{Subtemas y diapositivas (2/2)}\par\medskip
\begin{columns}[T,onlytextwidth]
\begin{column}{.48\textwidth}
\IndexEntry{nav-frame-014}{Aparición de $\mathcal R^{ab}$}
\IndexEntry{nav-frame-015}{Resultado del segundo cálculo}
\IndexEntry{nav-frame-016}{Comparación con el cálculo escalar}
\IndexEntry{nav-frame-017}{Consecuencia: simetría de \(\mathcal R^{ab}\)}
\IndexEntry{nav-frame-018}{Caso Einstein--Hilbert: cálculo de \(P^{abcd}\)}
\IndexEntry{nav-frame-019}{Cálculo de \(\mathcal R_{ab}\) para Einstein--Hilbert}
\end{column}
\begin{column}{.48\textwidth}
\IndexEntry{nav-frame-020}{Por qué \(\mathcal R_{ab}=R_{ab}\)}
\IndexEntry{nav-frame-021}{Constante cosmológica: momentos y tensor $\mathcal{R}_{ab}$}
\IndexEntry{nav-frame-022}{Casos de referencia: momentos y contracción}
\IndexEntry{nav-map-1-6}{Mapa de cierre}
\IndexEntry{nav-frame-023}{Resumen: identidad principal y caso Einstein--Hilbert}
\end{column}
\end{columns}
\vfill
\hyperlink{nav-index-1-1}{\scriptsize Anterior}\qquad
\hyperlink{nav-map-1-0}{\scriptsize Inicio de sección}
\end{frame}

\begin{frame}[label=nav-index-2-1]{Índice: Variación de la acción (LL)}
\footnotesize\raggedright
\textcolor{structure.fg}{Subtemas y diapositivas (1/2)}\par\medskip
\begin{columns}[T,onlytextwidth]
\begin{column}{.48\textwidth}
\IndexEntry{nav-map-2-0}{Mapa de apertura}
\IndexEntry{nav-frame-024}{Variación de la acción}
\IndexEntry{nav-frame-025}{Cambio de variable: \(g_{ab}\) a \(g^{ab}\)}
\IndexEntry{nav-frame-026}{Variación del lagrangiano}
\IndexEntry{nav-frame-027}{Variación de curvatura: separación métrica}
\IndexEntry{nav-frame-028}{Aplicación de la identidad de Palatini}
\IndexEntry{nav-frame-029}{Reemplazo de \(\delta\Gamma\) en \(P^{abcd}\delta R_{abcd}\)}
\end{column}
\begin{column}{.48\textwidth}
\IndexEntry{nav-frame-030}{Integración por partes: traslado de las derivadas}
\IndexEntry{nav-frame-031}{Integración por partes: volumen y frontera}
\IndexEntry{nav-frame-032}{Término de borde}
\IndexEntry{nav-frame-033}{Resultado variacional: tensor de campo y frontera}
\IndexEntry{nav-frame-034}{Identidad de Bianchi desde la covariancia}
\IndexEntry{nav-frame-035}{Identidad de Bianchi generalizada de Padmanabhan}
\IndexEntry{nav-frame-036}{Forma de la identidad en Lanczos--Lovelock}
\end{column}
\end{columns}
\vfill
\hyperlink{nav-index-2-2}{\scriptsize Más subtemas}\qquad
\hyperlink{nav-map-2-0}{\scriptsize Inicio de sección}
\end{frame}

\begin{frame}[label=nav-index-2-2]{Índice: Variación de la acción (LL)}
\footnotesize\raggedright
\textcolor{structure.fg}{Subtemas y diapositivas (2/2)}\par\medskip
\begin{columns}[T,onlytextwidth]
\begin{column}{.48\textwidth}
\IndexEntry{nav-frame-037}{Caso Einstein--Hilbert: Recuperación del tensor de Einstein}
\IndexEntry{nav-frame-038}{Término de frontera para Einstein--Hilbert}
\IndexEntry{nav-frame-039}{Constante cosmológica: ecuaciones y frontera}
\IndexEntry{nav-frame-040}{Casos de referencia: ecuaciones y frontera}
\IndexEntry{nav-frame-041}{Enlace: de la variación a la corriente de Noether}
\IndexEntry{nav-frame-042}{Difeomorfismos y corriente de Noether}
\IndexEntry{nav-frame-043}{Término de frontera bajo un difeomorfismo}
\end{column}
\begin{column}{.48\textwidth}
\IndexEntry{nav-frame-044}{Corriente de Noether explícita}
\IndexEntry{nav-frame-045}{Potencial de Noether: verificación de su divergencia}
\IndexEntry{nav-frame-046}{Einstein--Hilbert: corriente y potencial de Noether}
\IndexEntry{nav-frame-047}{Constante cosmológica: corriente y potencial de Noether}
\IndexEntry{nav-frame-048}{Casos de referencia: corriente y potencial}
\IndexEntry{nav-map-2-6}{Mapa de cierre}
\IndexEntry{nav-frame-049}{Resumen: ecuación métrica y término de frontera}
\end{column}
\end{columns}
\vfill
\hyperlink{nav-index-2-1}{\scriptsize Anterior}\qquad
\hyperlink{nav-map-2-0}{\scriptsize Inicio de sección}
\end{frame}

\begin{frame}[label=nav-index-3-1]{Índice: Condición de segundo orden}
\footnotesize\raggedright
\textcolor{structure.fg}{Subtemas y diapositivas (1/1)}\par\medskip
\begin{columns}[T,onlytextwidth]
\begin{column}{.48\textwidth}
\IndexEntry{nav-map-3-0}{Mapa de apertura}
\IndexEntry{nav-frame-051}{Cuarto orden y condición de divergencia nula de \(P\)}
\end{column}
\begin{column}{.48\textwidth}
\IndexEntry{nav-map-3-3}{Mapa de cierre}
\IndexEntry{nav-frame-052}{Resumen: condición de Lanczos--Lovelock}
\end{column}
\end{columns}
\vfill
\hyperlink{nav-map-3-0}{\scriptsize Inicio de sección}
\end{frame}

\begin{frame}[label=nav-index-4-1]{Índice: Estructura holográfica de la acción LL}
\footnotesize\raggedright
\textcolor{structure.fg}{Subtemas y diapositivas (1/1)}\par\medskip
\begin{columns}[T,onlytextwidth]
\begin{column}{.48\textwidth}
\IndexEntry{nav-map-4-0}{Mapa de apertura}
\IndexEntry{holo-ll-01}{Acción y datos de borde: construcción mecánica}
\IndexEntry{holo-ll-02}{Momento fijo, campos y fase de frontera}
\IndexEntry{holo-ll-03}{Separación geométrica: integración por partes}
\IndexEntry{holo-ll-04}{Einstein--Hilbert: volumen, superficie e identidad}
\end{column}
\begin{column}{.48\textwidth}
\IndexEntry{holo-ll-05}{Einstein--Hilbert: momento y variación de frontera}
\IndexEntry{holo-ll-06}{Lanczos--Lovelock: homogeneidad y Bianchi}
\IndexEntry{holo-ll-07}{Segundo gradiente: dos integraciones por partes}
\IndexEntry{holo-ll-08}{Homogeneidad y deducción de la identidad LL}
\IndexEntry{holo-ll-10}{Resumen LL: identidad, controles y alcance}
\end{column}
\end{columns}
\vfill
\hyperlink{nav-map-4-0}{\scriptsize Inicio de sección}
\end{frame}

\begin{frame}[label=nav-index-5-1]{Índice: Generalización con campo escalar}
\footnotesize\raggedright
\textcolor{structure.fg}{Subtemas y diapositivas (1/2)}\par\medskip
\begin{columns}[T,onlytextwidth]
\begin{column}{.48\textwidth}
\IndexEntry{nav-map-5-0}{Mapa de apertura}
\IndexEntry{nav-frame-072}{Generalización del lagrangiano}
\IndexEntry{nav-frame-073}{Diccionario de los momentos generalizados}
\IndexEntry{nav-frame-074}{Variación del lagrangiano generalizado}
\IndexEntry{nav-frame-075}{Variaciones: resultado gravitacional y aporte escalar}
\IndexEntry{nav-frame-076}{Identidad generalizada para \(\Rcal_{ab}\)}
\end{column}
\begin{column}{.48\textwidth}
\IndexEntry{nav-frame-077}{Variación completa de la acción}
\IndexEntry{nav-frame-078}{Ecuación métrica como extensión de $L(g,R)$}
\IndexEntry{nav-frame-079}{Variación: caso geométrico y extensión escalar}
\IndexEntry{nav-frame-080}{Bianchi con escalar: variación por difeomorfismos}
\IndexEntry{nav-frame-081}{Bianchi con escalar: identidad fuera de las soluciones}
\end{column}
\end{columns}
\vfill
\hyperlink{nav-index-5-2}{\scriptsize Más subtemas}\qquad
\hyperlink{nav-map-5-0}{\scriptsize Inicio de sección}
\end{frame}

\begin{frame}[label=nav-index-5-2]{Índice: Generalización con campo escalar}
\footnotesize\raggedright
\textcolor{structure.fg}{Subtemas y diapositivas (2/2)}\par\medskip
\begin{columns}[T,onlytextwidth]
\begin{column}{.48\textwidth}
\IndexEntry{nav-frame-082}{Corriente de Noether en \(L(g,R,\phi,\nabla\phi)\)}
\IndexEntry{nav-frame-083}{Bianchi generalizada y corriente conservada}
\IndexEntry{nav-frame-084}{Cancelación explícita del sector escalar}
\IndexEntry{nav-frame-085}{Potencial de Noether generalizado}
\IndexEntry{nav-frame-086}{Comparación con el formalismo de Padmanabhan}
\end{column}
\begin{column}{.48\textwidth}
\IndexEntry{nav-frame-087}{Forma útil para teorías con simetría de desplazamiento}
\IndexEntry{nav-frame-088}{Identidad de Bianchi en términos de los momentos}
\IndexEntry{nav-frame-089}{Ecuaciones de campo}
\IndexEntry{nav-map-5-6}{Mapa de cierre}
\IndexEntry{nav-frame-090}{Resumen escalar: dinámica, Bianchi y Noether}
\end{column}
\end{columns}
\vfill
\hyperlink{nav-index-5-1}{\scriptsize Anterior}\qquad
\hyperlink{nav-map-5-0}{\scriptsize Inicio de sección}
\end{frame}

\begin{frame}[label=nav-index-6-1]{Índice: Estudio de casos}
\footnotesize\raggedright
\textcolor{structure.fg}{Subtemas y diapositivas (1/2)}\par\medskip
\begin{columns}[T,onlytextwidth]
\begin{column}{.48\textwidth}
\IndexEntry{nav-map-6-0}{Mapa de apertura}
\IndexEntry{nav-frame-091}{Acción concreta: separación de los sectores}
\IndexEntry{nav-frame-092}{Sector cinético: momentos de curvatura y métrica}
\IndexEntry{nav-frame-093}{Momentos escalares de \(L_{\mathrm{cin}}\)}
\IndexEntry{nav-frame-094}{Variación directa de \(L_{\mathrm{cin}}\)}
\IndexEntry{nav-frame-095}{Integración por partes del término escalar}
\IndexEntry{nav-frame-096}{Reconstrucción del sector \(L_0\)}
\end{column}
\begin{column}{.48\textwidth}
\IndexEntry{nav-frame-097}{Resultado final para \(L_0\)}
\IndexEntry{nav-frame-098}{Sector de acoplamiento \(Q_\beta\)}
\IndexEntry{nav-frame-099}{Momentos de $Q_\beta$: cálculo y uso posterior}
\IndexEntry{nav-frame-100}{Momento de curvatura \(P_\beta^{abcd}\)}
\IndexEntry{nav-frame-101}{Momento métrico \(P^{(\beta)}_{ab}\)}
\IndexEntry{nav-frame-102}{Momento métrico: grupo proporcional a \(3\)}
\IndexEntry{nav-frame-103}{Momento métrico: grupo con signo negativo}
\end{column}
\end{columns}
\vfill
\hyperlink{nav-index-6-2}{\scriptsize Más subtemas}\qquad
\hyperlink{nav-map-6-0}{\scriptsize Inicio de sección}
\end{frame}

\begin{frame}[label=nav-index-6-2]{Índice: Estudio de casos}
\footnotesize\raggedright
\textcolor{structure.fg}{Subtemas y diapositivas (2/2)}\par\medskip
\begin{columns}[T,onlytextwidth]
\begin{column}{.48\textwidth}
\IndexEntry{nav-frame-104}{Momento métrico \(P^{(\beta)}_{ab}\): resultado}
\IndexEntry{nav-frame-105}{Momentos de $Q_\beta$: curvatura y métrica}
\IndexEntry{nav-frame-106}{Momento escalar \(\Pi_\beta^a\)}
\IndexEntry{nav-frame-107}{Momento escalar \(\Pi_\beta^a\): resultado}
\IndexEntry{nav-frame-108}{Momentos de $Q_\beta$: gradiente y campo escalar}
\IndexEntry{nav-frame-109}{Contracción \(\Rcal^{(\beta)}_{(ab)}\)}
\IndexEntry{nav-frame-110}{Término \(-2\nabla\nabla P_\beta\)}
\end{column}
\begin{column}{.48\textwidth}
\IndexEntry{nav-frame-111}{Variación de la acción \(I_{Q_\beta}\)}
\IndexEntry{nav-frame-112}{Ecuaciones asociadas a \(Q_\beta\)}
\IndexEntry{nav-frame-113}{Término de frontera de \(Q_\beta\)}
\IndexEntry{nav-frame-114}{Chequeo de Bianchi para $L_0+Q_\beta$}
\IndexEntry{nav-map-6-6}{Mapa de cierre}
\IndexEntry{nav-frame-115}{Resumen EQT: ecuaciones de la teoría completa}
\end{column}
\end{columns}
\vfill
\hyperlink{nav-index-6-1}{\scriptsize Anterior}\qquad
\hyperlink{nav-map-6-0}{\scriptsize Inicio de sección}
\end{frame}

\begin{frame}[label=nav-index-7-1]{Índice: Estructura holográfica generalizada}
\footnotesize\raggedright
\textcolor{structure.fg}{Subtemas y diapositivas (1/1)}\par\medskip
\begin{columns}[T,onlytextwidth]
\begin{column}{.48\textwidth}
\IndexEntry{nav-map-7-0}{Mapa de apertura}
\IndexEntry{holo-gen-01}{Hipótesis y deducción del defecto de escalamiento}
\IndexEntry{holo-gen-02}{Separación exacta y orden de las derivadas}
\IndexEntry{holo-gen-03}{Operador de escalamiento del volumen}
\IndexEntry{holo-gen-05}{Identidad volumen--superficie con defecto}
\end{column}
\begin{column}{.48\textwidth}
\IndexEntry{holo-gen-06}{Superficie efectiva: significado y límite}
\IndexEntry{holo-gen-07}{Sector $Q_\beta$: peso y contracción de $P_\beta$}
\IndexEntry{holo-gen-08}{Sector $Q_\beta$: divergencia y superficie efectiva}
\IndexEntry{holo-gen-09}{Chequeo sobre el ansatz estático}
\IndexEntry{holo-gen-10}{Resumen: identidad con defecto y frontera}
\end{column}
\end{columns}
\vfill
\hyperlink{nav-map-7-0}{\scriptsize Inicio de sección}
\end{frame}

\begin{frame}[label=nav-index-8-1]{Índice: Discusión de resultados}
\footnotesize\raggedright
\textcolor{structure.fg}{Subtemas y diapositivas (1/1)}\par\medskip
\begin{columns}[T,onlytextwidth]
\begin{column}{.48\textwidth}
\IndexEntry{nav-map-8-0}{Mapa de apertura}
\IndexEntry{nav-frame-116}{Condición de Lanczos--Lovelock: $\nabla_a P^{abcd}$}
\IndexEntry{nav-frame-117}{Divergencia de $P_\beta^{abcd}$}
\IndexEntry{nav-frame-118}{Expresión general para $\nabla_aP_\beta^{abcd}$}
\IndexEntry{nav-frame-119}{Evaluación sobre el ansatz de la solución}
\IndexEntry{nav-frame-120}{No se cumple la condición de Lanczos--Lovelock}
\end{column}
\begin{column}{.48\textwidth}
\IndexEntry{nav-frame-121}{Lectura radial de la componente calculada}
\IndexEntry{nav-map-8-4}{Mapa de cierre}
\IndexEntry{nav-frame-122}{Resumen del contraste: EQT no es LL en general}
\IndexEntry{nav-frame-123}{Síntesis del recorrido teórico}
\IndexEntry{nav-frame-124}{Referencias}
\end{column}
\end{columns}
\vfill
\hyperlink{nav-map-8-0}{\scriptsize Inicio de sección}
\end{frame}
